%% Generated by Sphinx.
\def\sphinxdocclass{report}
\documentclass[letterpaper,10pt,english]{sphinxmanual}
\ifdefined\pdfpxdimen
   \let\sphinxpxdimen\pdfpxdimen\else\newdimen\sphinxpxdimen
\fi \sphinxpxdimen=.75bp\relax
\ifdefined\pdfimageresolution
    \pdfimageresolution= \numexpr \dimexpr1in\relax/\sphinxpxdimen\relax
\fi
%% let collapsible pdf bookmarks panel have high depth per default
\PassOptionsToPackage{bookmarksdepth=5}{hyperref}

\PassOptionsToPackage{booktabs}{sphinx}
\PassOptionsToPackage{colorrows}{sphinx}

\PassOptionsToPackage{warn}{textcomp}
\usepackage[utf8]{inputenc}
\ifdefined\DeclareUnicodeCharacter
% support both utf8 and utf8x syntaxes
  \ifdefined\DeclareUnicodeCharacterAsOptional
    \def\sphinxDUC#1{\DeclareUnicodeCharacter{"#1}}
  \else
    \let\sphinxDUC\DeclareUnicodeCharacter
  \fi
  \sphinxDUC{00A0}{\nobreakspace}
  \sphinxDUC{2500}{\sphinxunichar{2500}}
  \sphinxDUC{2502}{\sphinxunichar{2502}}
  \sphinxDUC{2514}{\sphinxunichar{2514}}
  \sphinxDUC{251C}{\sphinxunichar{251C}}
  \sphinxDUC{2572}{\textbackslash}
\fi
\usepackage{cmap}
\usepackage[T1]{fontenc}
\usepackage{amsmath,amssymb,amstext}
\usepackage{babel}



\usepackage{tgtermes}
\usepackage{tgheros}
\renewcommand{\ttdefault}{txtt}



\usepackage[Bjarne]{fncychap}
\usepackage{sphinx}

\fvset{fontsize=auto}
\usepackage{geometry}


% Include hyperref last.
\usepackage{hyperref}
% Fix anchor placement for figures with captions.
\usepackage{hypcap}% it must be loaded after hyperref.
% Set up styles of URL: it should be placed after hyperref.
\urlstyle{same}

\addto\captionsenglish{\renewcommand{\contentsname}{Contents:}}

\usepackage{sphinxmessages}
\setcounter{tocdepth}{1}



\title{COMP206 Learning Journal}
\date{Jul 26, 2026}
\release{}
\author{Hector Barquero}
\newcommand{\sphinxlogo}{\vbox{}}
\renewcommand{\releasename}{}
\makeindex
\begin{document}

\ifdefined\shorthandoff
  \ifnum\catcode`\=\string=\active\shorthandoff{=}\fi
  \ifnum\catcode`\"=\active\shorthandoff{"}\fi
\fi

\pagestyle{empty}
\sphinxmaketitle
\pagestyle{plain}
\sphinxtableofcontents
\pagestyle{normal}
\phantomsection\label{\detokenize{index::doc}}


\sphinxstepscope


\chapter{Assignment four journal}
\label{\detokenize{AssignmentFour/main:assignment-four-journal}}\label{\detokenize{AssignmentFour/main::doc}}

\begin{savenotes}\sphinxattablestart
\sphinxthistablewithglobalstyle
\centering
\begin{tabulary}{\linewidth}[t]{TT}
\sphinxtoprule
\sphinxtableatstartofbodyhook
\sphinxAtStartPar
Time spent
&
\sphinxAtStartPar
Jan 2026 to Aug 2026
\\
\sphinxhline
\sphinxAtStartPar
Student name
&
\sphinxAtStartPar
Hector Barquero
\\
\sphinxbottomrule
\end{tabulary}
\sphinxtableafterendhook\par
\sphinxattableend\end{savenotes}

\begin{sphinxadmonition}{tip}{Tip:}
\sphinxAtStartPar
\sphinxstylestrong{How to use this journal}

\sphinxAtStartPar
Learning journal contains program design notes, changelogs for versions as I made updates, reflections, and reading notes for the units of each assignment. There’s also practice questions from the textbook which were recommended to be added to the learning journal for each assignment chapter.
\end{sphinxadmonition}


\section{Program design}
\label{\detokenize{AssignmentFour/main:program-design}}

\subsection{Concept and scope}
\label{\detokenize{AssignmentFour/main:concept-and-scope}}
\sphinxAtStartPar
This project is a compact parser\sphinxhyphen{}based adventure game inspired by Lewis
Carroll’s \sphinxstyleemphasis{Alice’s Adventures in Wonderland}. You (the player) represent Alice and
begin beside the rabbit hole. The White Rabbit has lost four important
objects: the Golden Key, Pocket Watch, Painted Rose and Queen’s Tart. You
must descend into Wonderland, explore its connected locations, collect the four
treasures and return them to the Hall of Doors. The Hall is the game’s safe
room. After all treasures are stashed there, Alice must use the DRINK ME Bottle
to become small enough to enter the tiny door and reach the Quiet Garden.

\sphinxAtStartPar
The game uses recognizable characters, places and objects while giving them a new game objective.
The game scope demonstrates object\sphinxhyphen{}oriented design, external data
loading, a command parser, state changes, inventory management, random events,
scoring, winning and losing. It remains small enough to test every location and
important command manually and I’ve included some instructions for fast tracking to
certain winning or losing states.


\subsection{Player experience and command style}
\label{\detokenize{AssignmentFour/main:player-experience-and-command-style}}
\sphinxAtStartPar
All communication is text. On the first visit to a location, the game prints a
long description. Later visits print a shorter reminder. It then lists visible
items, characters and available exits. The player enters short commands. Input
is converted to lowercase, punctuation becomes whitespace and repeated spaces
are removed. Therefore \sphinxcode{\sphinxupquote{GO NORTH!}}, \sphinxcode{\sphinxupquote{go north}} and \sphinxcode{\sphinxupquote{Go, north}} are treated
consistently.

\sphinxAtStartPar
Direction shortcuts are supported. \sphinxcode{\sphinxupquote{go north}}, \sphinxcode{\sphinxupquote{north}} and \sphinxcode{\sphinxupquote{n}} all
produce the same internal command. Other aliases are loaded from
\sphinxcode{\sphinxupquote{actions.txt}}. Examples include \sphinxcode{\sphinxupquote{inventory}}, \sphinxcode{\sphinxupquote{invent}} and \sphinxcode{\sphinxupquote{i}} or
\sphinxcode{\sphinxupquote{take}}, \sphinxcode{\sphinxupquote{get}} and \sphinxcode{\sphinxupquote{pick up}}. The parser separates an action from its
object and optional target. For example, \sphinxcode{\sphinxupquote{throw cup at hatter}} becomes the
action \sphinxcode{\sphinxupquote{throw}}, object \sphinxcode{\sphinxupquote{cup}} and target \sphinxcode{\sphinxupquote{hatter}}.

\sphinxAtStartPar
The opening message states the goal. \sphinxcode{\sphinxupquote{help}} lists valid commands and \sphinxcode{\sphinxupquote{score}}
shows score, turns, energy, courage and treasure progress. Invalid commands do
not end the game. They produce a direct message and allow another attempt.


\subsection{World plan}
\label{\detokenize{AssignmentFour/main:world-plan}}
\sphinxAtStartPar
The map contains eleven locations. The rabbit hole is intentionally one\sphinxhyphen{}way.
Once Alice moves down from the Riverbank, no exit leads back above ground.

\begin{sphinxVerbatim}[commandchars=\\\{\}]
\PYG{n}{Sister}\PYG{l+s+s1}{\PYGZsq{}}\PYG{l+s+s1}{s Lawn \PYGZlt{}\PYGZhy{}\PYGZgt{} Riverbank}
                      \PYG{o}{|}
                    \PYG{n}{down}
                      \PYG{n}{v}
                 \PYG{n}{Rabbit} \PYG{n}{Hole}
                      \PYG{o}{|}
                    \PYG{n}{down}
                      \PYG{n}{v}
              \PYG{n}{Hall} \PYG{n}{of} \PYG{n}{Doors} \PYG{p}{[}\PYG{n}{SAFE}\PYG{p}{]}
       \PYG{n}{north} \PYG{o}{/}      \PYG{n}{east} \PYG{o}{|} \PYG{n}{west}       \PYGZbs{} \PYG{n}{south}
            \PYG{n}{v}            \PYG{n}{v}              \PYG{n}{v}
     \PYG{n}{Pool} \PYG{n}{of} \PYG{n}{Tears}  \PYG{n}{Mushroom} \PYG{n}{Forest}  \PYG{n}{Mad} \PYG{n}{Tea} \PYG{n}{Party}
            \PYG{o}{|} \PYG{n}{east}       \PYG{o}{|} \PYG{n}{south}\PYG{o}{/}\PYG{n}{east}   \PYG{o}{|} \PYG{n}{south}
            \PYG{n}{v}            \PYG{n}{v}              \PYG{n}{v}
     \PYG{n}{Caucus} \PYG{n}{Garden} \PYG{o}{\PYGZlt{}}\PYG{o}{\PYGZhy{}}\PYG{o}{\PYGZgt{}} \PYG{n}{Queen}\PYG{l+s+s1}{\PYGZsq{}}\PYG{l+s+s1}{s Garden \PYGZlt{}\PYGZhy{}\PYGZgt{} Royal Kitchen}
                           \PYG{o}{|}
                   \PYG{k}{return} \PYG{n}{north} \PYG{n}{to} \PYG{n}{Hall}

\PYG{n}{Hall} \PYG{n}{of} \PYG{n}{Doors} \PYG{o}{\PYGZhy{}}\PYG{o}{\PYGZhy{}}\PYG{n}{down}\PYG{p}{,} \PYG{n}{when} \PYG{n}{unlocked} \PYG{o+ow}{and} \PYG{n}{Alice} \PYG{o+ow}{is} \PYG{n}{small}\PYG{o}{\PYGZhy{}}\PYG{o}{\PYGZhy{}}\PYG{o}{\PYGZgt{}} \PYG{n}{Quiet} \PYG{n}{Garden} \PYG{p}{[}\PYG{n}{WIN}\PYG{p}{]}
\end{sphinxVerbatim}

\sphinxAtStartPar
The Hall of Doors has five exits, satisfying the requirement for a room with
more than three exits. It connects the major areas and acts as a safe gathering
point. Stashed items do not count against the six\sphinxhyphen{}item carrying capacity. This
encourages the player to return rather than carrying every useful object at
once.

\sphinxAtStartPar
The locations are Sister’s Lawn, Riverbank, Rabbit Hole, Hall of Doors, Pool of
Tears, Caucus Garden, Mushroom Forest, Mad Tea Party, Queen’s Garden, Royal
Kitchen and Quiet Garden. Every location has an outgoing route and at least one
incoming connection. The Rabbit Hole prevents reverse travel by omitting an
upward exit. The Quiet Garden has an exit in the data model, although reaching
it immediately completes the game.


\subsection{Items and treasures}
\label{\detokenize{AssignmentFour/main:items-and-treasures}}
\sphinxAtStartPar
The game defines 16 interactive items, exceeding the required 12.
They include the White Rabbit’s Fan, Marmalade Jar, DRINK ME Bottle, EAT ME
Cake, Thimble, Box of Comfits, Golden Key, Mushroom Piece, Tea Cup, Pocket
Watch, Bread and Butter, Painted Rose, Croquet Flamingo, Hedgehog Croquet Ball,
Pepper Shaker and Queen’s Tart.

\sphinxAtStartPar
Items have an identifier, display name, aliases, description, current location,
takeable flag, treasure flag, point value, use message and state flags. Their
location changes to \sphinxcode{\sphinxupquote{inventory}} when taken, a room identifier when dropped,
\sphinxcode{\sphinxupquote{stash}} when stored safely or \sphinxcode{\sphinxupquote{consumed}} after single\sphinxhyphen{}use food and drink.
This produces one consistent representation instead of separate room and
inventory copies.

\sphinxAtStartPar
Food restores energy. The fan and painted rose protect Alice during hostile
encounters. The pocket watch reports the current turn. The flamingo and
hedgehog can be combined indirectly by carrying both and using either on the
croquet field. The DRINK ME Bottle sets the small\sphinxhyphen{}size state needed for the
final door. The EAT ME Cake and Mushroom Piece restore ordinary size.


\subsection{Characters and encounters}
\label{\detokenize{AssignmentFour/main:characters-and-encounters}}
\sphinxAtStartPar
Eight characters are defined: White Rabbit, Dodo, Caterpillar, Mad Hatter,
Dormouse, Cheshire Cat, Queen of Hearts and Duchess. Each stores aliases, a
location, description, friendly flag, dialogue and random encounter text.
Talking to a character prints location\sphinxhyphen{}specific dialogue and awards a small
first\sphinxhyphen{}conversation score.

\sphinxAtStartPar
Random encounters can occur after state\sphinxhyphen{}changing actions. A friendly encounter
adds atmosphere without punishment. A hostile Queen or Duchess encounter may
reduce courage. Carrying the tart makes a Queen encounter more likely. The fan
or painted rose grants temporary protection. The \sphinxcode{\sphinxupquote{throw}} action also has a
random outcome: the target either ducks or is bumped harmlessly. Random events
use a Mersenne Twister generator. A seed can be supplied for repeatable testing,
and random encounters can be disabled when testing a fixed route.


\subsection{Classes and responsibilities}
\label{\detokenize{AssignmentFour/main:classes-and-responsibilities}}
\sphinxAtStartPar
\sphinxcode{\sphinxupquote{Game}} is the highest class. It owns the repositories, current location,
Inventory, Control object, score, turn count, energy, courage and other game
flags. It loads data, runs the loop and dispatches parsed commands to methods.

\sphinxAtStartPar
\sphinxcode{\sphinxupquote{Location}} stores one room. \sphinxcode{\sphinxupquote{Locations}} loads and manages all Location
objects. Each Location contains descriptions, visited and safe flags, and a map
from directions to destination identifiers.

\sphinxAtStartPar
\sphinxcode{\sphinxupquote{Item}} stores one object. \sphinxcode{\sphinxupquote{Items}} loads all Item objects and finds them by
identifier, alias and current location. \sphinxcode{\sphinxupquote{Character}} and \sphinxcode{\sphinxupquote{Characters}} use
the same entity\sphinxhyphen{}and\sphinxhyphen{}repository pattern for non\sphinxhyphen{}player characters.

\sphinxAtStartPar
\sphinxcode{\sphinxupquote{Action}} stores a canonical action, aliases and help text. \sphinxcode{\sphinxupquote{Actions}} loads
the action vocabulary and finds the longest alias matching the beginning of a
command. Longest matching is necessary so \sphinxcode{\sphinxupquote{pick up}} is selected as a phrase
instead of treating only \sphinxcode{\sphinxupquote{pick}} as an action.

\sphinxAtStartPar
\sphinxcode{\sphinxupquote{Inventory}} maintains a vector of carried item identifiers and enforces a
capacity of six. It does not duplicate full Item objects. \sphinxcode{\sphinxupquote{Control}} reads a
line, normalizes it, recognizes direction shortcuts, matches a loaded Action
and separates the object and target. \sphinxcode{\sphinxupquote{TextUtils}} supplies reusable string
operations such as trimming, splitting, joining and command normalization.


\subsection{External text files}
\label{\detokenize{AssignmentFour/main:external-text-files}}
\sphinxAtStartPar
\sphinxcode{\sphinxupquote{locations.txt}} contains room names, long and short descriptions, the safe
flag and exits. \sphinxcode{\sphinxupquote{items.txt}} defines object properties and starting locations.
\sphinxcode{\sphinxupquote{characters.txt}} contains descriptions, dialogue and encounter text.
\sphinxcode{\sphinxupquote{actions.txt}} contains canonical command names, aliases and help descriptions.
Each non\sphinxhyphen{}comment record uses vertical bars between fields and commas inside
alias or exit lists. Loading fails with a specific message if a file is missing
or a record has the wrong number of fields. Cross\sphinxhyphen{}references are checked so an
item, character or exit cannot silently point to an unknown location.


\subsection{Game state, scoring and ending}
\label{\detokenize{AssignmentFour/main:game-state-scoring-and-ending}}
\sphinxAtStartPar
A successful action advances the turn. Energy decreases every third turn.
Consumable items restore it. If energy reaches zero, the game ends. Alice begins
with three courage points. Hostile interactions and some random encounters can
reduce courage; reaching zero also ends the game.

\sphinxAtStartPar
The player earns points for first visits, first conversations, selected items,
the croquet interaction and safely stashing treasures. The score is secondary
to the explicit win condition. The tiny door checks two facts: all four
required treasures currently have \sphinxcode{\sphinxupquote{stash}} as their location, and Alice has
used the DRINK ME Bottle. Passing both checks allows travel to the Quiet Garden,
where the game prints the final score and ends successfully.


\subsection{Algorithmic scale}
\label{\detokenize{AssignmentFour/main:algorithmic-scale}}
\sphinxAtStartPar
Command normalization is linear in command length, \sphinxcode{\sphinxupquote{O(c)}}. Action matching is
\sphinxcode{\sphinxupquote{O(a * s)}} for \sphinxcode{\sphinxupquote{a}} aliases and short command strings. Item and character
searches are linear in their small collections, \sphinxcode{\sphinxupquote{O(i)}} and \sphinxcode{\sphinxupquote{O(k)}}. Location
lookup and exit lookup use ordered maps, giving \sphinxcode{\sphinxupquote{O(log r)}} for \sphinxcode{\sphinxupquote{r}} rooms or
exits. The complete game loop is \sphinxcode{\sphinxupquote{O(t * p)}}, where \sphinxcode{\sphinxupquote{t}} is the number of
turns and \sphinxcode{\sphinxupquote{p}} is the bounded per\sphinxhyphen{}command parsing and object\sphinxhyphen{}search work. With
the fixed assignment data, each turn behaves as constant practical work.


\section{Program test plan}
\label{\detokenize{AssignmentFour/main:program-test-plan}}
\begin{sphinxVerbatim}[commandchars=\\\{\}]
\PYG{n}{g}\PYG{o}{+}\PYG{o}{+} \PYG{o}{\PYGZhy{}}\PYG{n}{o} \PYG{n}{Wonderland}\PYG{o}{.}\PYG{n}{exe} \PYG{o}{*}\PYG{o}{.}\PYG{n}{cpp}
\end{sphinxVerbatim}

\sphinxAtStartPar
Expected result: compilation completes without warnings or errors. Run with ./Wonderland.exe

\begin{sphinxVerbatim}[commandchars=\\\{\}]
\PYG{o}{.}\PYG{o}{/}\PYG{n}{Wonderland}\PYG{o}{.}\PYG{n}{exe} \PYG{o}{\PYGZhy{}}\PYG{o}{\PYGZhy{}}\PYG{n}{seed} \PYG{l+m+mi}{1} \PYG{o}{\PYGZhy{}}\PYG{o}{\PYGZhy{}}\PYG{n}{no}\PYG{o}{\PYGZhy{}}\PYG{n}{random} \PYG{o}{\PYGZlt{}} \PYG{n}{tests}\PYG{o}{/}\PYG{n}{win\PYGZus{}commands}\PYG{o}{.}\PYG{n}{txt}
\end{sphinxVerbatim}

\sphinxAtStartPar
Expected results:
\begin{itemize}
\item {} 
\sphinxAtStartPar
The program loads all four external data files.

\item {} 
\sphinxAtStartPar
The four treasures can be collected and stashed.

\item {} 
\sphinxAtStartPar
The output states that all four treasures are safe.

\item {} 
\sphinxAtStartPar
Using the bottle enables the final downward movement.

\item {} 
\sphinxAtStartPar
The Quiet Garden and \sphinxcode{\sphinxupquote{You win}} messages appear.

\end{itemize}

\begin{sphinxVerbatim}[commandchars=\\\{\}]
\PYG{o}{.}\PYG{o}{/}\PYG{n}{Wonderland}\PYG{o}{.}\PYG{n}{exe} \PYG{o}{\PYGZhy{}}\PYG{o}{\PYGZhy{}}\PYG{n}{seed} \PYG{l+m+mi}{1} \PYG{o}{\PYGZhy{}}\PYG{o}{\PYGZhy{}}\PYG{n}{no}\PYG{o}{\PYGZhy{}}\PYG{n}{random} \PYG{o}{\PYGZlt{}} \PYG{n}{tests}\PYG{o}{/}\PYG{n}{parser\PYGZus{}commands}\PYG{o}{.}\PYG{n}{txt}
\end{sphinxVerbatim}

\sphinxAtStartPar
Expected results:
\begin{itemize}
\item {} 
\sphinxAtStartPar
Punctuation and capitalization are ignored.

\item {} 
\sphinxAtStartPar
\sphinxcode{\sphinxupquote{invent}} displays the inventory.

\item {} 
\sphinxAtStartPar
An unavailable direction is rejected.

\item {} 
\sphinxAtStartPar
An unknown action is rejected without terminating the loop.

\item {} 
\sphinxAtStartPar
Taking and dropping the fan changes inventory state.

\end{itemize}

\begin{sphinxVerbatim}[commandchars=\\\{\}]
\PYG{o}{.}\PYG{o}{/}\PYG{n}{Wonderland}\PYG{o}{.}\PYG{n}{exe} \PYG{o}{\PYGZhy{}}\PYG{o}{\PYGZhy{}}\PYG{n}{data} \PYG{n}{tests}\PYG{o}{/}\PYG{n}{missing}\PYG{o}{\PYGZhy{}}\PYG{n}{data}
\end{sphinxVerbatim}

\sphinxAtStartPar
Expected result: the program reports the missing \sphinxcode{\sphinxupquote{locations.txt}} file and
returns a failure status.
\begin{enumerate}
\sphinxsetlistlabels{\arabic}{enumi}{enumii}{}{.}%
\item {} 
\sphinxAtStartPar
Fill all six inventory slots and attempt to take another item. Confirm that
the capacity message appears.

\item {} 
\sphinxAtStartPar
Attempt \sphinxcode{\sphinxupquote{stash key}} outside the Hall of Doors. Confirm that its rejected.

\item {} 
\sphinxAtStartPar
Stash an item in the Hall, then take it again. Confirm that it returns to
inventory.

\item {} 
\sphinxAtStartPar
Attempt to enter the tiny door without four treasures. Confirm its locked.

\item {} 
\sphinxAtStartPar
Stash all treasures but do not use the bottle. Confirm Alice is too large.

\item {} 
\sphinxAtStartPar
Use the cake after the bottle. Confirm Alice is again too large for the door.

\item {} 
\sphinxAtStartPar
Use \sphinxcode{\sphinxupquote{talk to queen}} without protection until courage reaches zero. Confirm
the loss condition.

\item {} 
\sphinxAtStartPar
Use food while energy is below its maximum. Confirm recovery does not exceed
30.

\item {} 
\sphinxAtStartPar
Carry the tart with random encounters enabled and a fixed seed. Confirm a
Queen encounter can occur.

\item {} 
\sphinxAtStartPar
Change an item location to an unknown ID in \sphinxcode{\sphinxupquote{items.txt}}. Confirm startup
validation reports the item ID.

\end{enumerate}


\section{Overview and reflections}
\label{\detokenize{AssignmentFour/main:overview-and-reflections}}
\sphinxAtStartPar
I began by reviewing the history and interaction style of \sphinxstyleemphasis{Colossal Cave
Adventure}. The research showed that the game established the familiar pattern
of a narrator describing a location, a player entering short commands, and the
program changing a model of rooms, objects and progress. Sources describing the
game emphasize movement between named locations, contextual object commands,
inventory, treasure collection and a score. I also noted the early parser’s
separation between motion words and action words. In particular, historical
versions treated commands such as \sphinxcode{\sphinxupquote{go north}} and \sphinxcode{\sphinxupquote{north}} as equivalent.
That directly influenced the command handling in this project.

\sphinxAtStartPar
The Ubuntu manual page for Open Adventure was useful because it describes the
parser behaviour from a player’s perspective. It identifies the historic
motion and action distinction and explains that motion words are handled
immediately. I did not reproduce that parser exactly because some of its odd
interpretations would be confusing in a new student project. Instead, I kept
the useful conventions: short commands, action\sphinxhyphen{}object structure, direction
abbreviations and forgiving capitalization.

\sphinxAtStartPar
I read the Project Gutenberg edition of \sphinxstyleemphasis{Alice’s Adventures in Wonderland} to
select locations, characters and objects. The final design uses the riverbank,
rabbit hole, hall of doors, pool of tears, caucus race, mushroom scene, tea
party, Queen’s garden and kitchen. Objects such as the fan, marmalade jar,
labelled bottle and cake, golden key, mushroom, watch, painted rose, flamingo and
tarts gave the game enough items to support inventory actions without inventing
an unrelated fantasy setting.

\sphinxAtStartPar
The research also suggested separating program logic from descriptive data.
Historical Adventure implementations contained substantial location,
vocabulary and message data in addition to executable logic. I used that idea
in a smaller object\sphinxhyphen{}oriented form: four text files define the world while C++
classes implement the rules. This also satisfies the assignment requirement
that actions, items, characters and location descriptions be loaded at startup.

\sphinxAtStartPar
The first decision was to avoid adapting the entire book. That would create too
many rooms and special cases to test reliably. I selected a connected group of
recognizable scenes and gave them a new, explicit objective. The White Rabbit
has lost four treasures. Alice must recover and store them, then enter the tiny
door. The objective is stated before the first command so the player doesnt
need prior knowledge of the novel.

\sphinxAtStartPar
I chose the Hall of Doors as the safe room because its central to the selected
scenes and naturally supports many doors. It has five exits, so it also satisfies
the requirement that one room have more than three exits. The map branches from
this hub into the Pool of Tears, Mushroom Forest, Mad Tea Party and Queen’s
Garden. Cross\sphinxhyphen{}connections reduce unnecessary backtracking and make the world
feel less like a list of isolated rooms.

\sphinxAtStartPar
Inventory capacity is six. This is not required for basic storage, but it gives
the safe room a practical purpose. Alice can carry all four treasures and one or
two tools, but collecting every object at once is impossible. The player must
choose, drop or stash. Items are represented by one object whose \sphinxcode{\sphinxupquote{location}}
field changes. I selected this model after considering separate lists in each
room, the inventory and the safe area. Separate lists create synchronization
problems because an item can accidentally exist in two places. A single location
value makes the invariant clear: each item is in exactly one place.

\sphinxAtStartPar
The parser is deliberately limited rather than pretending to understand free
English. It normalizes text and searches the loaded action aliases. The longest
matching alias is selected, which permits a phrase such as \sphinxcode{\sphinxupquote{pick up}}. The
remaining words become the object, except that \sphinxcode{\sphinxupquote{at}} and \sphinxcode{\sphinxupquote{on}} can divide an
object from a target. This supports commands such as \sphinxcode{\sphinxupquote{throw cup at hatter}}.
Direction words are handled separately so \sphinxcode{\sphinxupquote{go north}}, \sphinxcode{\sphinxupquote{north}} and \sphinxcode{\sphinxupquote{n}} map
to the same command.

\sphinxAtStartPar
I used first\sphinxhyphen{}visit and repeat descriptions to match a common text\sphinxhyphen{}adventure
pattern and the assignment requirements. The Location object stores both forms
and a visited flag. The first visit gives two exploration points and prints the
long description. Later visits print the shorter version. The rule is contained
in one method rather than repeated in every movement branch.

\sphinxAtStartPar
Random encounters were added only after deterministic movement and item logic
worked. Random behaviour can make testing difficult, so the program supports a
fixed seed and a \sphinxcode{\sphinxupquote{\sphinxhyphen{}\sphinxhyphen{}no\sphinxhyphen{}random}} option. Normal play still includes encounters.
The Queen becomes more likely when Alice carries the tart, demonstrating that
inventory can influence an encounter. Hostile encounters reduce courage only
sometimes, and protective items can prevent the loss. This keeps randomness
meaningful without making a winning route depend entirely on luck.

\sphinxAtStartPar
The largest structural challenge was deciding how the required plural classes
should interact. I created small entity structures, such as Location and Item,
plus collection classes named Locations and Items. The collections are
responsible for loading and searching, while Game is responsible for game
rules. This prevents Game from parsing every file itself and keeps the exact
required component names visible in the design.

\sphinxAtStartPar
External data parsing introduced several possible errors: missing files, wrong
field counts, invalid numeric values and references to locations that do not
exist. Initially it would have been possible for an item to load with a bad
room identifier and simply never appear. I added validation after loading. Game
checks the starting room, every item location, every character location and
every exit destination. A failed check prints the responsible identifier and
ends startup.

\sphinxAtStartPar
A specific identifier bug appeared during the first winning\sphinxhyphen{}route test. The
same function used to normalize player commands was also used on data IDs.
Because command normalization replaces punctuation with spaces, an identifier
such as \sphinxcode{\sphinxupquote{hall\_of\_doors}} became \sphinxcode{\sphinxupquote{hall of doors}}. Hard\sphinxhyphen{}coded state checks still
used the underscore form, so the final\sphinxhyphen{}door restrictions were bypassed and the
win location was not recognized. I corrected the design by treating internal
identifiers differently from commands. IDs are trimmed and converted to
lowercase but retain underscores. Player aliases are still normalized
aggressively. This test demonstrated why machine identifiers and user\sphinxhyphen{}facing
text should have separate rules.

\sphinxAtStartPar
Another parser issue involved \sphinxcode{\sphinxupquote{use object with target}}. The separators
\textasciigrave{}\textasciigrave{} on \textasciigrave{}\textasciigrave{} and \textasciigrave{}\textasciigrave{} with \textasciigrave{}\textasciigrave{} have different lengths. The first implementation used
the same substring offset for both, leaving part of the word \sphinxcode{\sphinxupquote{with}} in the
target. I changed the parser to record which separator matched and remove the
correct number of characters.

\sphinxAtStartPar
C++ object ownership was kept simple. Game owns repository objects directly and
owns Control through \sphinxcode{\sphinxupquote{std::unique\_ptr}} because Control cannot be constructed
until Actions has loaded. Items and characters are stored as values in maps.
Inventory stores only IDs, so it cannot own or duplicate an Item. Methods obtain
pointers to repository values for short operations. No pointer is stored after
the map changes, and the maps are not structurally modified after loading.

\sphinxAtStartPar
Compilation testing used C++17 with \sphinxcode{\sphinxupquote{\sphinxhyphen{}Wall}}, \sphinxcode{\sphinxupquote{\sphinxhyphen{}Wextra}} and \sphinxcode{\sphinxupquote{\sphinxhyphen{}pedantic}}.
The complete project compiled with no warnings. The first automated scenario
uses a command transcript and disables random encounters. It travels from the
riverbank through the rabbit hole, collects the four treasures, returns to the
Hall of Doors, stashes them, uses the bottle and enters the Quiet Garden. The
expected output includes the four\sphinxhyphen{}treasure message, the win message and a final
score.

\sphinxAtStartPar
Parser tests include capitalization, punctuation and aliases. \sphinxcode{\sphinxupquote{LOOK!!!}} must
behave as \sphinxcode{\sphinxupquote{look}}. \sphinxcode{\sphinxupquote{INVENT}} must behave as \sphinxcode{\sphinxupquote{inventory}}. Movement tests use
both full and one\sphinxhyphen{}letter forms. Unknown actions must report an error without
changing location. Invalid movement must also preserve the current location.

\sphinxAtStartPar
Inventory testing verifies take, list and drop. A separate manual sequence fills
all six slots and attempts to take a seventh item. The expected response explains
that something must be dropped or stashed. Safe\sphinxhyphen{}room testing verifies that
\sphinxcode{\sphinxupquote{stash}} fails elsewhere, succeeds in the Hall, removes the item from carried
inventory and makes it visible in the safe\sphinxhyphen{}room list. Taking a stashed item back
out reduces the current treasure count.

\sphinxAtStartPar
Item tests cover size and energy state. The DRINK ME Bottle is consumed once and
sets the small state. The EAT ME Cake and Mushroom Piece restore normal size.
Food increases energy without passing the maximum. The pocket watch reports the
turn count. Carrying both croquet objects and using one in the Queen’s Garden
awards the croquet score only once.

\sphinxAtStartPar
Character tests call \sphinxcode{\sphinxupquote{talk}} for friendly and hostile characters. A first
conversation awards points, while repeated conversation doesnt repeatedly add
that score. Hostile conversation reduces courage when protection is inactive.
The fan or rose supplies temporary protection. Throw tests use a fixed seed to
confirm both the parser and the item’s transfer from inventory back to the room.

\sphinxAtStartPar
Failure testing renames or omits a data file and checks for a non\sphinxhyphen{}zero exit. I
also changed an exit destination to an invalid identifier and confirmed that
startup rejected it. These tests matter because the external text files are
part of the final program, not optional supporting material.

\sphinxAtStartPar
The random system was tested in two modes. With a fixed seed, the same sequence
of actions produces repeatable encounters. With a random\sphinxhyphen{}device seed, repeated
short games produce different encounter timing. Randomness is intentionally not
used for core puzzle success. A player can always win through correct movement,
inventory and size state if energy and courage are managed.

\sphinxAtStartPar
The parser doesnt understand unrestricted sentences or pronouns. It expects a
known action followed by a recognizable object name. This is appropriate for the
course scope and is made explicit by the help command. The data format also doesnt
support vertical bars inside descriptive text because that character is the
field separator.

\sphinxAtStartPar
The final design meets the required minimums and provides a complete game loop.
It includes eleven locations, eight characters, sixteen items, external data,
random events, a safe room, inventory changes, long and short descriptions,
multiple command forms, loss conditions and a defined win. The major design
lesson was that a text adventure is easiest to manage when world data, parsing
and state\sphinxhyphen{}changing rules have separate responsibilities.


\section{Research notes}
\label{\detokenize{AssignmentFour/main:research-notes}}

\subsection{Research and game concept}
\label{\detokenize{AssignmentFour/main:research-and-game-concept}}\begin{itemize}
\item {} 
\sphinxAtStartPar
Researched the short\sphinxhyphen{}command parser, room exploration, inventory and treasure\sphinxhyphen{}scoring conventions used by Colossal Cave Adventure. (v0.1\sphinxhyphen{}v0.2)

\item {} 
\sphinxAtStartPar
Reviewed Alice’s Adventures in Wonderland for locations, characters and interactive objects suitable for a smaller game. (v0.1\sphinxhyphen{}v0.2)

\item {} 
\sphinxAtStartPar
Defined the goal as recovering four lost treasures, returning them to a safe room and entering the tiny door. (v0.3)

\item {} 
\sphinxAtStartPar
Reduced the adaptation to a connected set of scenes that could be completely documented and tested. (v0.3\sphinxhyphen{}v0.4)

\end{itemize}


\subsection{World map and external data}
\label{\detokenize{AssignmentFour/main:world-map-and-external-data}}\begin{itemize}
\item {} 
\sphinxAtStartPar
Designed eleven connected locations with the Rabbit Hole as a one\sphinxhyphen{}way transition into Wonderland. (v0.4)

\item {} 
\sphinxAtStartPar
Made the Hall of Doors a safe room with five exits and a stash for treasure storage. (v0.4\sphinxhyphen{}v0.5)

\item {} 
\sphinxAtStartPar
Added long first\sphinxhyphen{}visit descriptions and short repeat\sphinxhyphen{}visit descriptions to locations.txt. (v0.5)

\item {} 
\sphinxAtStartPar
Added sixteen items to items.txt, including four treasures and several usable food, size and protection items. (v0.5\sphinxhyphen{}v0.6)

\item {} 
\sphinxAtStartPar
Added eight characters with descriptions, dialogue, friendly states and random encounter text to characters.txt. (v0.6)

\item {} 
\sphinxAtStartPar
Added canonical actions, command aliases and help descriptions to actions.txt. (v0.6)

\item {} 
\sphinxAtStartPar
Added startup validation for missing files, malformed records and invalid location references. (v1.1\sphinxhyphen{}v1.2)

\end{itemize}


\subsection{Object\sphinxhyphen{}oriented class design}
\label{\detokenize{AssignmentFour/main:object-oriented-class-design}}\begin{itemize}
\item {} 
\sphinxAtStartPar
Created Game as the highest class and owner of the complete game state and execution loop. (v0.7)

\item {} 
\sphinxAtStartPar
Created Location and Locations classes to store, load and search the room map. (v0.7)

\item {} 
\sphinxAtStartPar
Created Item and Items classes to maintain one location state for every object. (v0.7\sphinxhyphen{}v0.8)

\item {} 
\sphinxAtStartPar
Created Character and Characters classes for fixed characters and random encounters. (v0.8)

\item {} 
\sphinxAtStartPar
Created Action and Actions classes to load valid commands and aliases. (v0.8)

\item {} 
\sphinxAtStartPar
Created Inventory with a six\sphinxhyphen{}item capacity so the safe room has a practical purpose. (v0.8\sphinxhyphen{}v0.9)

\item {} 
\sphinxAtStartPar
Created Control to read, normalize and parse player input before Game executes an action. (v0.9)

\end{itemize}


\subsection{Command parser and game logic}
\label{\detokenize{AssignmentFour/main:command-parser-and-game-logic}}\begin{itemize}
\item {} 
\sphinxAtStartPar
Added lowercase conversion, punctuation removal and whitespace normalization for all player commands. (v0.9)

\item {} 
\sphinxAtStartPar
Added support for go north, north and n as equivalent movement commands. (v0.9\sphinxhyphen{}v1.0)

\item {} 
\sphinxAtStartPar
Used longest\sphinxhyphen{}alias matching so multiword actions such as pick up are recognized correctly. (v1.0)

\item {} 
\sphinxAtStartPar
Added object and target parsing for commands such as throw cup at hatter and use fan on queen. (v1.0)

\item {} 
\sphinxAtStartPar
Added movement, look, take, drop, stash, use, talk, examine, throw, score, help and quit actions. (v1.0\sphinxhyphen{}v1.1)

\item {} 
\sphinxAtStartPar
Added first\sphinxhyphen{}visit scoring, inventory capacity, energy, courage and safe\sphinxhyphen{}room treasure tracking. (v1.1)

\item {} 
\sphinxAtStartPar
Added random character encounters and increased the chance of meeting the Queen while carrying the tart. (v1.1)

\item {} 
\sphinxAtStartPar
Added a fixed random seed and no\sphinxhyphen{}random mode for repeatable scripted testing. (v1.1)

\end{itemize}


\subsection{Testing and corrections}
\label{\detokenize{AssignmentFour/main:testing-and-corrections}}\begin{itemize}
\item {} 
\sphinxAtStartPar
Compiled all source files with C++17, Wall, Wextra and pedantic warning flags. (v1.2)

\item {} 
\sphinxAtStartPar
Created a deterministic command transcript that collects and stashes all four treasures and reaches the winning room. (v1.2)

\item {} 
\sphinxAtStartPar
Found that command normalization changed underscore\sphinxhyphen{}based internal IDs into spaces and bypassed final\sphinxhyphen{}door state checks. (v1.2)

\item {} 
\sphinxAtStartPar
Separated internal identifier handling from player\sphinxhyphen{}command normalization so IDs retain their underscores. (v1.2\sphinxhyphen{}v1.3)

\item {} 
\sphinxAtStartPar
Corrected target extraction for use object with target because with and on use different separator lengths. (v1.3)

\item {} 
\sphinxAtStartPar
Verified the final door remains locked without four stashed treasures and rejects Alice when she is not small enough. (v1.3)

\item {} 
\sphinxAtStartPar
Verified the complete winning route ends in the Quiet Garden and prints the final score and turn count. (v1.3)

\item {} 
\sphinxAtStartPar
Added parser, missing\sphinxhyphen{}data, inventory\sphinxhyphen{}capacity, hostile\sphinxhyphen{}character and cross\sphinxhyphen{}reference validation cases to the test plan. (v1.3\sphinxhyphen{}v1.4)

\end{itemize}


\section{Sources and references}
\label{\detokenize{AssignmentFour/main:sources-and-references}}\begin{enumerate}
\sphinxsetlistlabels{\arabic}{enumi}{enumii}{}{.}%
\item {} 
\sphinxAtStartPar
cppreference.com. (n.d.). C++ reference. Retrieved Jan 2026 \sphinxhyphen{} Aug 2026, from \sphinxurl{https://cppreference.com/cpp}

\item {} 
\sphinxAtStartPar
Braunschweig, D., \& Busbee, K. L. (2018). Programming fundamentals: A modular structured approach (2nd ed.). Rebus Community. \sphinxurl{https://press.rebus.community/programmingfundamentals/}

\item {} 
\sphinxAtStartPar
Hansen, J. A. (2013). The Rook’s guide to C++. Rook’s Guide Press. \sphinxurl{https://rooksguide.org/}

\item {} 
\sphinxAtStartPar
Downey, A. B. (1999). Think C++ (Version 1.1.0). Green Tea Press. \sphinxurl{https://www.greenteapress.com/thinkcpp/}

\item {} 
\sphinxAtStartPar
leetcode.com. (n.d.). \sphinxstyleemphasis{LeetCode}. Retrieved Jan 2025 \sphinxhyphen{} Aug 2026, from \sphinxurl{https://leetcode.com/}

\item {} 
\sphinxAtStartPar
Project Gutenberg, \sphinxstyleemphasis{Alice’s Adventures in Wonderland}:

\end{enumerate}
\begin{quote}

\sphinxAtStartPar
\sphinxurl{https://www.gutenberg.org/ebooks/11}
\end{quote}
\begin{enumerate}
\sphinxsetlistlabels{\arabic}{enumi}{enumii}{}{.}%
\item {} 
\sphinxAtStartPar
Ubuntu Manpages, \sphinxcode{\sphinxupquote{advent}} and the Open Adventure parser description:

\end{enumerate}
\begin{quote}

\sphinxAtStartPar
\sphinxurl{https://manpages.ubuntu.com/manpages/bionic/man6/advent.6.html}
\end{quote}
\begin{enumerate}
\sphinxsetlistlabels{\arabic}{enumi}{enumii}{}{.}%
\item {} 
\sphinxAtStartPar
Eric S. Raymond, \sphinxstyleemphasis{A Brief History of Colossal Cave Adventure}:

\end{enumerate}
\begin{quote}

\sphinxAtStartPar
\sphinxurl{https://www.catb.org/~esr/open-adventure/history.html}
\end{quote}
\begin{enumerate}
\sphinxsetlistlabels{\arabic}{enumi}{enumii}{}{.}%
\item {} 
\sphinxAtStartPar
Wikipedia, \sphinxstyleemphasis{Colossal Cave Adventure}, used for a broad gameplay and history (the assignment overview before confirming parser details with the manual page): \sphinxurl{https://en.wikipedia.org/wiki/Colossal\_Cave\_Adventure}

\end{enumerate}



\renewcommand{\indexname}{Index}
\printindex
\end{document}